%-------------------------------------------------------------------------------
\section{Background}
%-------------------------------------------------------------------------------

% ### 2. Background and Threat Model
In this section, we first introduce the typical ecosystem architecture and payment scenarios of third-party online payments. 
Then, we abstract two main payment models, and finally introduce the formal modeling tool Tamarin that we use.


% #### 2.1 Payment Ecosystem Overview
\subsection{Third-Party Online Payment Ecosystem}

Currently, there is no unified industry standard for third-party online payment systems, 
as different payment providers maintain their own proprietary protocols and implementations. 
Through extensive review of official documentation and developer guides\cite{WeChat2025APIV3,Alipay2025OpenPlatform}, 
combined with practical analysis of payment workflows,
we have summarized the typical ecosystem architecture of third-party online payments.
A typical third-party online payment ecosystem involves four main participants:

\begin{figure}[t]
    \centering
    \includegraphics[width=\linewidth]{img/black.png}
    \caption{Typical Third-Party Online Payment Process}
    \label{fig:third_party_payment}
\end{figure}

% - 参与方:用户 / 商家 / 第三方支付商
\begin{enumerate} 
  \item \textit{\textbf{User (U):}}
    An individual who initiates the payment for goods or services.
  \item \textit{\textbf{Merchant Frontend (MF):}}
    The interface through which the user interacts to select goods and initiate payment,
    typically a website or mobile application.
  \item \textit{\textbf{Merchant Server (MS):}}
    The server-side component of the merchant that processes orders and communicates with the payment provider.
  \item \textit{\textbf{Third-Party Payment System (PS):}}
    A service provider that handles the actual payment transactions (e.g., Alipay and WeChat Pay).
    The \emph{PS} comprises a payment application (\emph{PS-App}) and a payment server (\emph{PS-Server}).
    Because the interaction details between these components are proprietary and often carried over private secure channels,
    we treat the application and server together as a single black-box system.
\end{enumerate}


% - 支付流程的典型阶段
As illustrated in Figure.~\ref{fig:third_party_payment}, we abstract a typical third-party payment process into the following phases:
\begin{enumerate}
  \item[\rom{1}.] \textit{\textbf{Order Phase:}}
    \emph{U} purchases products and selects a payment provider on \emph{MF}. 
    The \emph{MF} sends the product and payment provider information to \emph{MS}.
  \item[\rom{2}.] \textit{\textbf{Prepayment Phase:}} 
    Based on the received information, \emph{MS} generates the corresponding payment parameters and returns them to \emph{MF}. 
    \emph{MF} forwards the payment parameters to the corresponding \emph{PS}.
  \item[\rom{3}.] \textit{\textbf{Payment Phase:}} 
    The \emph{PS} guides \emph{U} through the completion of the payment operation.
  \item[\rom{4}.] \textit{\textbf{Notification Phase:}}  
    After the payment is completed, the \emph{PS} sends a payment result notification to the \emph{MS}. 
    Meanwhile, the \emph{PS} also presents the payment result page to \emph{U}.
\end{enumerate}

It is important to note that, during the \emph{Prepayment} phase, there are two primary approaches for \emph{MS} to generate payment parameters.
We classify these approaches into two payment models, which will be discussed in detail in Section~\ref{sec:payment-models}.
    

% #### 2.2 Payment Scenarios
\subsection{Payment Scenarios}

\emph{PS} obtains payment information from \emph{MF} via \emph{PS-App} on the \emph{U}'s mobile device.
However, due to the diverse forms of \emph{MF}, the interaction methods between \emph{PS-App} and the \emph{MF} also vary. 
As previously outlined, \emph{MF} manifests in various forms, including native mobile applications, mobile web pages, and desktop web pages. 
These distinct forms differentiate corresponding payment scenarios, 
for which payment providers have designed specific interaction methods between the \emph{MF} and the payment system (\emph{PS}). 
In the following, we introduce three common payment scenarios:

\begin{figure}[t]
    \centering
    \includegraphics[width=\linewidth]{img/black.png}
    \caption{Payment Scenario \Rom{1} --- In-app Payment}
    \label{fig:inapp_payment}
\end{figure}

\begin{figure}[t]
    \centering
    \includegraphics[width=\linewidth]{img/black.png}
    \caption{Payment Scenario \Rom{2} --- Mobile Web Payment}
    \label{fig:mobile_web_payment}
\end{figure}

\begin{figure}[t]
    \centering
    \includegraphics[width=\linewidth]{img/black.png}
    \caption{Payment Scenario \Rom{3} --- QR Code Payment}
    \label{fig:qr_code_payment}
\end{figure}

% 3. App 内跳转支付
\paragraph{In-app Payment.}
This payment scenario takes place within native mobile applications.
After the user selects a product, the merchant application invokes the \emph{PS-App} through the Software Development Kit (SDK) provided by the payment service provider,
transmits the required payment parameters, and subsequently completes the payment workflow.

% 2. 手机浏览器跳转支付
\paragraph{Mobile Web Payment.}
In mobile web payment, the user interacts with the merchant through a mobile browser rather than a native application.
Similar to in-app payment, the merchant’s mobile web page initiates a payment request to the \emph{PS-App} using a URL scheme.

% 1. PC 网站扫码支付
\paragraph{QR Code Payment.}
QR code payment differs fundamentally from the previous two scenarios in terms of device interaction.
In this case, the user browses the merchant’s website using a desktop browser, while the payment is executed on a mobile device.
Since the merchant’s web page cannot directly redirect to a mobile payment application, the payment service provider employs a QR code as a bridging mechanism between the PC and mobile platforms.
After the user selects a product, the merchant website generates a payment QR code, which is scanned by the \emph{PS-App} to obtain the payment parameters and complete the transaction.
 
% #### 2.3 Payment Models
\subsection{Payment Models}\label{sec:payment-models}

As discussed earlier, during the \emph{Prepayment} phase, protocol designs vary across different providers and scenarios.
Here, we introduce two payment models adopted by two payment providers, WeChat Pay and Alipay, across three payment scenarios.

\begin{figure}[t]
  \centering
  \includegraphics[width=\linewidth]{img/black.png}
  \caption{Payment Model A}
  \label{fig:payment_model_a}
\end{figure}

\begin{figure}[t]
  \centering
  \includegraphics[width=\linewidth]{img/black.png}
  \caption{Payment Model B}
  \label{fig:payment_model_b}
\end{figure}

\begin{enumerate} 
  \item \textit{\textbf{Payment Model A:}}
    In this model, the payment parameters are generated directly by \emph{MS}.
    Figure.~\ref{fig:payment_model_a} illustrates in detail the process by which \emph{MS} generates the order parameters.
    Based on the $goods$ selected by \emph{U}, \emph{MS} calculates the payable $amount$ and the order $description$, and generates a unique order identifier~($id_{order}$).
    \emph{MS} then signs the order parameters using the merchant's private key $Sk_{M}$, producing a signature value $sign_M$.
    During the subsequent $Payment$ phase, the \emph{PS} verifies the signature value to ensure that the order parameters have not been tampered with.
    Payment Model A is adopted by Alipay in the In-app Payment and Mobile Web Payment scenarios.
  \item \textit{\textbf{Payment Model B:}}
    In contrast to Model~A, in Model~B, the payment parameters are generated by \emph{PS}.
    Similarly, the merchant is required to sign the critical order parameters.
    However, in this model, the order parameters and the signature value cannot be directly used as payment parameters. Instead, they are sent to \emph{PS} as a prepayment request.
    At this point, the \emph{PS} verifies the merchant's signature value to ensure the integrity of the order parameters.
    If the verification succeeds, the \emph{PS} generates a payment identifier~($id_{prepay}$) and returns it to the merchant as the payment parameter.
    During the subsequent $Payment$ phase, the \emph{PS} identifies the order information based on this random value and completes the payment.
    Payment Model~B is adopted by Alipay in the QR Code Payment scenario and by WeChat Pay across all three payment scenarios.
\end{enumerate}


\subsection{Tamarin Prover}

Tamarin Prover is a symbolic verification framework widely used for analyzing real-world security protocols, 
including EMV~\cite{Radu2022EMVRelayProtection,Basin2021EMVStandard,Basin2021CardBrandMixup}, TLS~\cite{Cremers2016AutomatedTLS,Cremers2017ComprehensiveTLS,Lee2019maTLS}, Bluetooth~\cite{Shi2023FormalBLEPairing,Jangid2023ExtrapolatingBluetoothPasskey}, and 5G-AKA~\cite{basin2018formal,cremers2019component}.

In Tamarin, protocols and adversaries are modeled as multiset-rewriting rules over a global state of facts. 
Cryptographic primitives are represented as abstract function symbols with algebraic properties defined by equations. 
This semantics induces a labeled transition system where executions correspond to traces of observable actions, enabling the specification of rich security properties (e.g., correspondence between user actions and payment outcomes) as first-order logic formulas over traces.

Tamarin supports unbounded sessions and explicitly models a Dolev-Yao adversary, allowing analysis under arbitrary concurrent sessions and adversarial control over selected channels. 
Crucially for payment protocols, it provides native support for mutable state and persistent facts, facilitating the representation of long-lived objects like order identifiers, transaction records, and order-account mappings. When a property fails, Tamarin either constructs a proof or generates a concrete, interpretable counterexample trace.

In this work, we employ Tamarin to model six deployed third-party payment flows under realistic threat models (e.g., compromised insecure channels). 
Its combination of multiset rewriting, persistent state, trace-based properties, and counterexample generation makes it particularly suited for analyzing the subtle relationships between orders, payments, and authorizations in online payment protocols.